Katharevousa Greek is under-served by modern Natural Language Processing (NLP) tools despite its importance in historical, legal, and parliamentary archives. We present a reproducible end-to-end workflow for constructing and evaluating a dependency parsing stack for Katharevousa from noisy OCR-derived sources. The pipeline combines text reconstruction, schema-constrained LLM-assisted annotation, deterministic CoNLL-U freezing, and fixed-split benchmark evaluation. The final gold snapshot contains 1,697 sentences. We compare feature-based, transformer, mBERT, and Stanza baselines under a shared test protocol. The best release candidate is an XLM-R parser (\texttt{v1\_opt}) with UPOS 0.8893, DEPREL F1 0.7250, UAS 0.6098, and LAS 0.5162. We further report custom Stanza training and negative tuning results to document practical trade-offs in low-resource historical NLP. The study contributes both a model benchmark and an operational methodology for reproducible language technology development on difficult archival corpora.
